% ============================================================
% 纯答案册·课时07-圆的方程 —— 工具/答案抽册器.py 抽册生成件（勿手改；第三档＝纯答案单独成册）
% 源件（只读）：C:/提示词/工作区/M3-第2章量产0913/成卷/练习件/课时07-圆的方程/main.tex（md5 c99db9764efb955af839b374d9ae142c）
%% 口径：题面不重印；逐题＝题号(印面号同正册)＋[答案]＋[详解]，值/详解照源件逐字
%       （钉值门硬断言：册值≡正册件内值≡值台账值tex，逐字全等）。册序＝正册装配序＝印面号 1..16 升序（同序同号）；键序断言基准＝题面库冻结 manifest canonical 键序。
%% 三档导出：整体 main-true／纯题 main-false（在产）｜本册＝纯答案册（本件）
%% 双档：ansbook-true.tex＝详解印本；ansbook-false.tex＝纯值速查本（详解不印）
% 编译：xelatex <壳> ×2，cwd＝本目录；sty＝qp-m3.sty 本地挂载（md5 7c3930362be8a0a2bdf21bbf8ac16573，锁校验）
% ============================================================
\documentclass[fontset=none]{ctexart}
\usepackage{qp-m3}
% —— 印面逐字符号路由补钉（承源件；− 走 TNR、▱ NSC 子块；禁改 sty 保 md5） ——
\xeCJKDeclareCharClass{Default}{"2212}%
\setCJKfamilyfont{nscblk}[Path=C:/提示词/工作区/字替对照-0909/variantF/fonts/,BoldFont=NSC-w600.ttf]{NSC-w500.ttf}%
\xeCJKDeclareSubCJKBlock{nscblk}{"25B1}%
\newcommand{\pxparallelogram}{{\CJKfamily{nscblk}▱}}%
% —— 断行微调（承母版实测档，三零门承重；\emergencystretch 不另设） ——
\xeCJKsetup{CJKglue={\hskip 0.10em plus 0.08em minus 0.05em}}%
% —— 纯答案册全线括线（长详解可跨栏断，承练习件全线括线口径；灰底盒不可跨栏断） ——
\ansblockgrayfalse
% —— 双档开关（false 档＝纯值速查：答案印、详解不印；件面开关，不动 sty 本体） ——
\newif\ifansbookdetail
\ifdefined\ansbookdetail
  \ifnum\ansbookdetail=0 \ansbookdetailfalse\else\ansbookdetailtrue\fi
\else
  \ansbookdetailtrue
\fi
% —— 页脚件名（承源件章名；件名＝<件型>答案册） ——
\renewcommand{\qpzhangming}{第二章\quad 平面解析几何}%
\renewcommand{\qpjianming}{练习件答案册}%

\begin{document}
\keshi{第7课时\quad 圆的方程}
\par\glueguard{1}\addvspace{2pt}\noindent\biaoqian{【纯答案册】}{\fangsong 题面不重印\quad 印面号与正册同号同序}\par\addvspace{2pt}

\begin{multicols}{2}
\emergencystretch=1em

\begin{ansblock}[2章-练-课时07-T1]
% ans:2章-练-课时07-T1
\ansitem{1}{D}
\ifansbookdetail
\ansnote{详解}{\((5a+1-1)^{2}+(12a)^{2}=25a^{2}+144a^{2}=169a^{2}<1\Longrightarrow |a|<1/13\)，选D。}
\fi
\end{ansblock}

\begin{ansblock}[2章-练-课时07-T2]
% ans:2章-练-课时07-T2
\ansitem{2}{D}
\ifansbookdetail
\ansnote{详解}{\((4a-1+1)^{2}+(3a+2-2)^{2}=16a^{2}+9a^{2}=25a^{2}\leqslant 25\Longrightarrow a^{2}\leqslant 1\Longrightarrow -1\leqslant a\leqslant 1\)，选D。}
\fi
\end{ansblock}

\begin{ansblock}[2章-练-课时07-T3]
% ans:2章-练-课时07-T3
\ansitem{3}{C}
\ifansbookdetail
\ansnote{详解}{表圆需\(m>0\)。原点在圆外\(\Longleftrightarrow (0-3)^{2}+(0+4)^{2}=25>m\)。故\(0<m<25\)，选C。}
\fi
\end{ansblock}

\begin{ansblock}[2章-练-课时07-T4]
% ans:2章-练-课时07-T4
\ansitem{4}{B}
\ifansbookdetail
\ansnote{详解}{\(x^{2}\)与\(y^{2}\)系数均为1。\ \(D^{2}+E^{2}-4F=4a^{2}+16a^{2}+40a=20a(a+2)>0\Longrightarrow a>0\)或\(a<-2\)，选B。}
\fi
\end{ansblock}

\begin{ansblock}[2章-练-课时07-T5]
% ans:2章-练-课时07-T5
\ansitem{5}{A}
\ifansbookdetail
\ansnote{详解}{表圆：\(m^{2}+(-2)^{2}-4\cdot 2>0\Longrightarrow m^{2}>4\Longrightarrow m>2\)或\(m<-2\)。点在圆外：\(1+4+m-4+2=m+3>0\Longrightarrow m>-3\)。取交：\((-3,-2)\cup(2,+\infty)\)，选A。}
\fi
\end{ansblock}

\begin{ansblock}[2章-练-课时07-T6]
% ans:2章-练-课时07-T6
\ansitem{6}{C}
\ifansbookdetail
\ansnote{详解}{过原点\(\Longleftrightarrow F=2m^{2}-6m+4=0\Longrightarrow m=1\)或2。\par\noindent 表圆条件：\(D^{2}+E^{2}-4F=4(m-1)^{2}+4(m-1)^{2}-4(2m^{2}-6m+4)=8(m^{2}-2m+1)-8m^{2}+24m-16=8m-8>0\Longrightarrow m>1\)。舍\(m=1\)，得\(m=2\)，选C。}
\fi
\end{ansblock}

\begin{ansblock}[2章-练-课时07-T7]
% ans:2章-练-课时07-T7
\ansitem{7}{ABC}
\ifansbookdetail
\ansnote{详解}{圆心\((2,0)\)。圆关于圆心中心对称（A对）；圆心在\(x\)轴上\(\Longrightarrow\)关于\(y=0\)对称（B对）；\(x+3y-2=0\)过\((2,0)\)（\(2+0-2=0\)）\(\Longrightarrow\)过圆心的直线是圆的对称轴（C对）；\(x-y+2=0\)：\(2-0+2=4\neq 0\)不过圆心，D错。选ABC。}
\fi
\end{ansblock}

\begin{ansblock}[2章-练-课时07-T8]
% ans:2章-练-课时07-T8
\ansitem{8}{D}
\ifansbookdetail
\ansnote{详解}{配方：\((x+1)^{2}+(y-m)^{2}=m^{2}+4m+5=(m+2)^{2}+1\)。\(r_{\min}=1\)当\(m=-2\)，此时圆心\((-1,-2)\)。\(|OC|=\sqrt{5}>r\)，圆上点到\(O\)距离最大\(=|OC|+r=\sqrt{5}+1\)，选D。}
\fi
\end{ansblock}

\begin{ansblock}[2章-练-课时07-T9]
% ans:2章-练-课时07-T9
\ansitem{9}{A}
\ifansbookdetail
\ansnote{详解}{圆心\((3,-5)\)。\(d=|12+15-2|/5=25/5=5\)。圆上点到直线最短距离\(=d-r\)（\(d>r\)时）\(=1\Longrightarrow r=4\)，选A。}
\fi
\end{ansblock}

\begin{ansblock}[2章-练-课时07-T10]
% ans:2章-练-课时07-T10
\ansitem{10}{CD}
\ifansbookdetail
\ansnote{详解}{圆心\((1,-2)\)。\(l\)平分圆\(\Longleftrightarrow l\)过圆心。截距相等分两类：①截距均为0：\(l\)过原点，\(k=-2/1=-2\)，\(l\)：\(2x+y=0\)（D对）；②截距\(a\neq 0\)：\(x/a+y/a=1\)，代入\((1,-2)\)：\((1-2)/a=1\Longrightarrow a=-1\)，\(l\)：\(x+y+1=0\)（C对）。\par\noindent 检验：A过\((1,-2)\)? \(1+2+1=4\neq 0\)；B过\((1,-2)\)? \(1+4=5\neq 0\)，均不过圆心。选CD。}
\fi
\end{ansblock}

\begin{ansblock}[2章-练-课时07-T11]
% ans:2章-练-课时07-T11
\ansitem{11}{x²+y²−4x+6y+8=0}
\ifansbookdetail
\ansnote{详解}{圆心在\(AB\)中垂线\(y=-3\)上，又在\(2x-y-7=0\)上：\(2x+3-7=0\Longrightarrow x=2\)，圆心\((2,-3)\)。\(r^{2}=4+(-3+4)^{2}=5\)。标准方程\((x-2)^{2}+(y+3)^{2}=5\)，展开：\(x^{2}+y^{2}-4x+6y+8=0\)。\par\noindent （验：\(A(0,-4)\)：\(16-24+8=0\)；\(B(0,-2)\)：\(4-12+8=0\)。）}
\fi
\end{ansblock}

\begin{ansblock}[2章-练-课时07-T12]
% ans:2章-练-课时07-T12
\ansitem{12}{(x−1)²+y²=9（答案不唯一）}
\ifansbookdetail
\ansnote{详解}{圆心\((a,0)\)，\(r^{2}=(1-a)^{2}+9\)。取\(a=1\)：\(r=3\)，得\((x-1)^{2}+y^{2}=9\)。（一般形式\((x-a)^{2}+y^{2}=(1-a)^{2}+9\)，\(a\in R\)均合。）}
\fi
\end{ansblock}

\begin{ansblock}[2章-练-课时07-T13]
% ans:2章-练-课时07-T13
\ansitem{13}{64；4；[3−2√2, 3+2√2]}
\ifansbookdetail
\ansnote{详解}{\((m+2)^{2}+(n+1)^{2}=|P-(-2,-1)|^{2}\)。圆心\(C(2,2)\)，\(r=3\)；\(|C-(-2,-1)|=\sqrt{16+9}=5\)。距离范围\([5-3,5+3]=[2,8]\)，平方范围\([4,64]\)，最大64、最小4。\par\noindent \(\sqrt{m^{2}+n^{2}}=|PO|\)，\(|OC|=2\sqrt{2}<3\)，故范围\([3-2\sqrt{2},3+2\sqrt{2}]\)。}
\fi
\end{ansblock}

\begin{ansblock}[2章-练-课时07-T14]
% ans:2章-练-课时07-T14
\ansitem{14}{(−2,−4)；5}
\ifansbookdetail
\ansnote{详解}{表圆需\(a^{2}=a+2\Longrightarrow a=-1\)或2。\(a=2\)：\(4x^{2}+4y^{2}+4x+8y+10=0\)，除以4配方得半径平方为负，舍。\(a=-1\)：\(x^{2}+y^{2}+4x+8y-5=0\)，\((x+2)^{2}+(y+4)^{2}=25\)。圆心\((-2,-4)\)，半径5。}
\fi
\end{ansblock}

\begin{ansblock}[2章-练-课时07-T15]
% ans:2章-练-课时07-T15
\ansitem{15}{(1) a≠−1；(2) 恒过 (0,0) 与 (16/5, −8/5)；(3) a=1/4}
\ifansbookdetail
\ansnote{详解}{(1) \(a\neq -1\)时\(x^{2}\)、\(y^{2}\)系数同为\(1+a\neq 0\)，两边除以\((1+a)\)化一般式，\(D^{2}+E^{2}-4F=(16+64a^{2})/(1+a)^{2}>0\)恒成立，表圆；\(a=-1\)时二次项消失，方程退化为直线\(-4x-8y=0\)，不合。故\(a\neq -1\)。\par\noindent (2) 方程按\(a\)整理：\((x^{2}+y^{2}-4x)+a(x^{2}+y^{2}+8y)=0\)。两定点须对一切\(a\)成立\(\Longleftrightarrow x^{2}+y^{2}-4x=0\)且\(x^{2}+y^{2}+8y=0\)。相减：\(4x+8y=0\)，\(x=-2y\)。代入\(x^{2}+y^{2}-4x=0\)：\(5y^{2}=0\)得\((0,0)\)；代入\(x^{2}+y^{2}+8y=0\)：\(5y^{2}+8y=0\Longrightarrow y=0\)（已得）或\(y=-8/5\)，\(x=16/5\)。定点\((0,0)\)与\((16/5,-8/5)\)。（验\((16/5,-8/5)\)：\(x^{2}+y^{2}=256/25+64/25=320/25\)；\(4x=64/5=320/25\)；\(8y=-64/5\)，两式皆零。）\par\noindent (3) 表圆时\(r^{2}=[4+16a^{2}]/(1+a)^{2}\)。令\(f(a)=(4+16a^{2})/(1+a)^{2}\)（\(a\neq -1\)），\(f'(a)=[32a(1+a)-2(4+16a^{2})]/(1+a)^{3}=(32a-8)/(1+a)^{3}\)，驻点\(a=1/4\)。符号按分支定号：①\(a>1/4\)：分子\(32a-8>0\)、分母\((1+a)^{3}>0\)，\(f'>0\)；②\(-1<a<1/4\)：分子\(<0\)、分母\(>0\)，\(f'<0\)；③\(a<-1\)：分子\(<0\)，但分母\((1+a)^{3}<0\)，\(f'>0\)——该支\(f\)单调递增且\(a\to -1^{-}\)时\(f\to+\infty\)、\(f\)恒大于\(16>16/5\)，不产生最小值。故\(f\)仅在\(a=1/4\)处取得最小值（\(16/5\)）。（几何：定点弦\((0,0)\)—\((16/5,-8/5)\)为公共弦，圆族中过两定点的圆以该弦为直径时半径最小，与上互证。）故\(a=1/4\)。}
\fi
\end{ansblock}

\begin{ansblock}[2章-练-课时07-T16]
% ans:2章-练-课时07-T16
\ansitem{16}{x²+y²+2x−4y+3=0}
\ifansbookdetail
\ansnote{详解}{圆心\((-D/2,-E/2)\)在\(x+y-1=0\)上：\(-D/2-E/2-1=0\Longrightarrow D+E=-2\)。\par\noindent \(r^{2}=(D^{2}+E^{2}-4F)/4=(D^{2}+E^{2}-12)/4=2\Longrightarrow D^{2}+E^{2}=20\)。由\(E=-2-D\)：\(D^{2}+(D+2)^{2}=20\Longrightarrow 2D^{2}+4D-16=0\Longrightarrow D^{2}+2D-8=0\Longrightarrow D=2\)或\(-4\)。\(D=2\Longrightarrow E=-4\)，圆心\((-1,2)\)第二象限；\(D=-4\Longrightarrow E=2\)，圆心\((2,-1)\)第四象限，舍。故\(D=2\)，\(E=-4\)，圆的一般方程\(x^{2}+y^{2}+2x-4y+3=0\)。（验\(r^{2}=(4+16-12)/4=2\)。）}
\fi
\end{ansblock}

% ---- 尾块（\tailfill[30mm] 定高参——钉档 §三.3：无参在平衡栏呈「内容尾随框」，贴栏底须定高参；实测无参致末页平衡 Overfull\vbox 12.99pt，定高 30mm 三零） ----
\tailfill[30mm]

\end{multicols}
\end{document}
